\chapter{PENUTUP}
\label{chap:penutup}

\section{Kesimpulan}
\label{sec:kesimpulan}

Berdasarkan hasil perancangan, implementasi, dan pengujian yang telah dilakukan, diperoleh beberapa kesimpulan sebagai berikut.

\begin{enumerate}

\item Sistem navigasi robot beroda otonom berhasil dikembangkan dengan mengintegrasikan LiDAR 2D dan kamera termal melalui pendekatan \textit{decision-level fusion} untuk menghasilkan nilai \textit{effective traversability} sebagai representasi kondisi lingkungan.

\item Logika fuzzy Mamdani berhasil diterapkan untuk menyesuaikan parameter navigasi APF secara adaptif berdasarkan jarak obstacle dan nilai \textit{effective traversability}, sehingga perilaku navigasi dapat disesuaikan dengan tingkat risiko lingkungan yang dihadapi robot.

\item Integrasi kamera termal dan logika fuzzy Mamdani terbukti meningkatkan keselamatan navigasi dibandingkan konfigurasi baseline. Peningkatan tersebut ditunjukkan oleh kenaikan \textit{Success Rate}, \textit{Collision-Free Rate}, serta \textit{minimum clearance}, terutama pada skenario yang mengandung obstacle berprofil rendah dan obstacle dinamis.

\item Metode yang diusulkan menghasilkan peningkatan waktu tempuh dan panjang lintasan dibandingkan konfigurasi baseline. Meskipun demikian, peningkatan tersebut masih berada dalam batas yang dapat diterima jika dibandingkan dengan peningkatan margin keselamatan yang diperoleh.

\end{enumerate}

\section{Saran}
\label{chap:saran}

Berdasarkan hasil penelitian yang telah diperoleh, beberapa pengembangan yang dapat dilakukan pada penelitian selanjutnya adalah sebagai berikut.

\begin{enumerate}

\item Melakukan implementasi dan pengujian sistem pada robot fisik untuk mengevaluasi pengaruh faktor lingkungan nyata, seperti variasi temperatur, ketidakpastian sensor, serta dinamika pergerakan obstacle yang tidak sepenuhnya dapat direpresentasikan pada lingkungan simulasi.

\item Mengembangkan metode pembentukan traversability termal yang lebih adaptif dengan memanfaatkan teknik segmentasi berbasis pembelajaran mendalam atau pendekatan per\-sepsi multimodal sehingga kualitas representasi lingkungan dapat ditingkatkan pada kondisi yang lebih kompleks.

\item Mengkaji metode fusi sensor yang lebih lanjut, seperti \textit{feature-level fusion} atau pendekatan probabilistik, untuk meningkatkan kualitas integrasi informasi antara LiDAR dan kamera termal sehingga sistem memiliki robustitas yang lebih baik terhadap ketidakpastian pengukuran sensor.

\item Mengintegrasikan metode navigasi lokal yang lebih robust terhadap permasalahan \textit{local minima} pada Artificial Potential Field, serta melakukan pengujian pada lingkungan dengan tingkat kompleksitas yang lebih tinggi, seperti jumlah obstacle yang lebih banyak, pergerakan obstacle yang lebih dinamis, atau skenario luar ruangan.

\end{enumerate}